\documentclass[11pt]{article}
\usepackage[utf8]{inputenc}
\usepackage[margin=1in]{geometry}
\usepackage{amsmath,amssymb}
\usepackage{booktabs}
\usepackage{hyperref}
\usepackage{natbib}
\usepackage{setspace}
\usepackage{float}
\usepackage{enumitem}
\usepackage{listings}
\onehalfspacing

\hypersetup{colorlinks=true,linkcolor=blue!70!black,citecolor=blue!70!black,urlcolor=blue!70!black}

\lstset{basicstyle=\ttfamily\small,breaklines=true,frame=single,columns=fullflexible}

\title{\textbf{ChronoFund: A Point-in-Time Fundamental\\Data Engine for Bias-Free Equity Research}}
\author{Tanishk Yadav\\[2pt]\textit{MS Financial Engineering}\\NYU Tandon School of Engineering\\New York, NY, USA\\[2pt]ORCID: 0009-0006-2382-9411}
\date{}

\begin{document}
\maketitle

\begin{abstract}
\noindent
Lookahead bias in fundamental data is a pervasive and underappreciated source of error in quantitative equity research. When a backtest uses a company's fiscal year 2023 financials on January 1, 2024---before the 10-K is filed with the SEC---it incorporates information unavailable to any real-world investor. This paper presents ChronoFund, a production-grade fundamental data engine that enforces point-in-time (PIT) correctness at four architectural layers: the filings index gates on SEC \texttt{acceptance\_datetime}, the filing selector raises \texttt{CutoffViolationError} on any temporal violation, the XBRL parser filters individual data points by cutoff date, and the Bloomberg mapper applies identical constraints. The engine parses raw XBRL from SEC EDGAR, maps US-GAAP and custom taxonomies, handles fiscal year misalignments, and outputs a unified schema spanning company metadata, filings, financial statements, and derived metrics. The system comprises 5,114 lines of Python across 30 modules with 756 lines of tests across 65 test cases. ChronoFund is packaged as an installable Python library with a CLI interface and serves as the data layer for the Regime-Aware Factor Backtest.

\medskip
\noindent\textbf{Keywords:} point-in-time data, lookahead bias, SEC EDGAR, XBRL, fundamental data, data engineering
\end{abstract}

\section{Introduction}

A persistent and often invisible source of error in quantitative equity backtests is lookahead bias in fundamental data. A company filing its 10-K for fiscal year ending December 31, 2023, typically does not have the filing accepted by the SEC until February or March 2024. Any backtest that uses these financial figures on or before the acceptance date is incorporating information that was not publicly available, producing artificially inflated performance metrics.

This problem is well-documented in academic literature \citep{banz1986}, yet the majority of open-source backtesting frameworks and commercial data providers either ignore it entirely or implement only rudimentary ``lag'' adjustments (e.g., adding 90 days to the fiscal period end). Such fixed lags are unreliable: actual filing delays vary from 30 to 180+ days depending on company size, filing type, and amendment history.

ChronoFund addresses this problem with a multi-layer architecture that makes temporal violations architecturally impossible, not merely discouraged.

\section{Architecture}

The system is organized as a Python package with four primary modules:

\subsection{EDGAR Module}
Handles all interaction with the SEC's Electronic Data Gathering, Analysis, and Retrieval system:
\begin{itemize}[nosep]
    \item \texttt{cik\_map}: Resolves ticker symbols to Central Index Key (CIK) numbers via the SEC company tickers endpoint.
    \item \texttt{client}: HTTP client with rate limiting (10 requests/second per SEC guidelines) and retry logic.
    \item \texttt{filings\_index}: Queries the EDGAR full-text search and submissions APIs, filtering filings by \texttt{acceptance\_datetime} $\leq$ cutoff date. This is the \textbf{primary temporal gate}.
    \item \texttt{xbrl/}: XBRL parsing pipeline comprising context extraction, fact fetching, taxonomy mapping, and fact selection.
\end{itemize}

\subsection{Bloomberg Module}
Ingests exported Bloomberg data in PDF and XLSX formats:
\begin{itemize}[nosep]
    \item \texttt{xlsx\_generic}: Parses generic Bloomberg Excel exports with column-per-period layout.
    \item \texttt{statement\_analysis\_pdf}: Extracts financial statements from Bloomberg FA (Financial Analysis) PDF exports.
    \item \texttt{segments\_pdf}: Parses business segment breakdown PDFs.
    \item \texttt{mapping}: Applies the same temporal cutoff logic as the EDGAR module.
\end{itemize}

\subsection{Snapshot Module}
Orchestrates the data assembly process:
\begin{itemize}[nosep]
    \item \texttt{builder}: Constructs point-in-time fundamental snapshots for a given universe of tickers and cutoff date.
    \item \texttt{selector}: Selects the most appropriate filing for each company-period pair, enforcing temporal constraints. Raises \texttt{CutoffViolationError} if any filing slips past the gate.
    \item \texttt{coverage}: Reports data completeness across the universe.
\end{itemize}

\subsection{Data Module}
Defines the unified output schema and validation rules:
\begin{itemize}[nosep]
    \item \texttt{schema}: Dataclass definitions for all output tables (company\_master, filings, statements\_income, statements\_balance, statements\_cashflow, derived\_metrics).
    \item \texttt{validation}: Schema validation, type checking, and referential integrity checks.
    \item \texttt{outputs}: Serialization to Parquet format.
\end{itemize}

\section{Point-in-Time Enforcement}

The PIT constraint is enforced at four independent layers:

\begin{table}[H]
\centering
\begin{tabular}{lp{7cm}}
\toprule
\textbf{Layer} & \textbf{Mechanism} \\
\midrule
\texttt{FilingsIndex} & Primary gate: \texttt{acceptance\_datetime} $\leq$ end of cutoff day \\
\texttt{FilingSelector} & Secondary assertion: raises \texttt{CutoffViolationError} if any filing slips through \\
\texttt{XBRLParser} & Passes \texttt{cutoff\_date} to every \texttt{select\_best\_fact\_for\_period()} call to exclude facts filed after cutoff \\
\texttt{BloombergMapper} & Skips columns whose \texttt{period\_end} $>$ \texttt{cutoff\_date} \\
\bottomrule
\end{tabular}
\end{table}

The use of \texttt{acceptance\_datetime}---the precise second the SEC accepted the filing---rather than \texttt{filing\_date} or \texttt{period\_end} is critical. The acceptance timestamp is the earliest moment at which any market participant could legally access the filing.

\section{XBRL Parsing}

Parsing XBRL is non-trivial due to the complexity of the XBRL taxonomy and the inconsistency of issuer implementations. ChronoFund's XBRL pipeline handles:

\begin{itemize}[nosep]
    \item \textbf{Context resolution}: Each XBRL fact is associated with a context specifying the entity, period (instant or duration), and optional dimensional qualifiers. The parser resolves contexts to determine which facts represent consolidated (vs. segment) data and which period they cover.
    \item \textbf{Taxonomy mapping}: Maps both US-GAAP standard concepts and company-specific extension concepts to a unified schema. For example, \texttt{us-gaap:RevenueFromContractWithCustomerExcludingAssessedTax} and \texttt{us-gaap:Revenues} both map to the ``Revenue'' field.
    \item \textbf{Fact selection}: When multiple facts exist for the same concept and period (common in amended filings), the parser selects the most recent fact filed before the cutoff date.
    \item \textbf{Fiscal year alignment}: Handles non-calendar fiscal years (e.g., Apple's September fiscal year end) by mapping fiscal periods to calendar quarters.
\end{itemize}

\section{Testing}

The test suite comprises 65 test cases across 5 test files (756 lines), covering:

\begin{itemize}[nosep]
    \item \textbf{Cutoff logic tests}: Verify that filings with \texttt{acceptance\_datetime} after the cutoff are excluded, that \texttt{CutoffViolationError} is raised on violations, and that the most recent pre-cutoff filing is selected.
    \item \textbf{CIK mapping tests}: Verify ticker-to-CIK resolution.
    \item \textbf{XBRL mapping tests}: Verify taxonomy concept mapping across US-GAAP versions.
    \item \textbf{Schema validation tests}: Verify output schema compliance and type correctness.
    \item \textbf{Bloomberg parser tests}: Verify PDF and XLSX parsing with known test fixtures.
\end{itemize}

\section{Limitations}

\begin{itemize}[nosep]
    \item \textbf{Coverage}: Limited to companies with clean XBRL filings. A surprising number of large-cap issuers submit malformed XBRL that fails parsing.
    \item \textbf{Bloomberg format dependency}: Bloomberg PDF/XLSX parsers rely on specific export format layouts that change without notice across Bloomberg terminal versions.
    \item \textbf{US-only}: No support for international exchanges, IFRS reporting, or non-US regulatory filings.
    \item \textbf{No real-time updates}: The engine processes batch snapshots, not streaming SEC filings.
    \item \textbf{Amended filings}: While the parser handles amendments (10-K/A), the amendment detection heuristic may miss some edge cases where only a subset of financial items are restated.
\end{itemize}

\section{Conclusion}

ChronoFund provides a rigorous, architecturally enforced solution to the point-in-time data problem in quantitative equity research. By gating data availability on SEC \texttt{acceptance\_datetime} at four independent layers, the engine makes temporal violations impossible rather than merely unlikely. The system is production-grade: 5,114 lines of modular Python code, 65 tests, proper package structure, and a CLI interface. It serves as the fundamental data layer for the Regime-Aware Factor Backtest, ensuring that all alpha signals in that strategy are derived exclusively from publicly available information.

\medskip
\noindent\textbf{Code Availability:} \url{https://github.com/tanishhky/chronofund-fundamental-engine}

\begin{thebibliography}{5}
\bibitem{banz1986} R.~W. Banz and W.~J. Breen, ``Sample-dependent results using accounting and market data: Some evidence,'' \textit{The Journal of Finance}, vol.~41, no.~4, pp.~779--793, 1986.
\bibitem{stambaugh1999} R.~F. Stambaugh, ``Predictive regressions,'' \textit{Journal of Financial Economics}, vol.~54, no.~3, pp.~375--421, 1999.
\bibitem{fama1993} E.~F. Fama and K.~R. French, ``Common risk factors in the returns on stocks and bonds,'' \textit{Journal of Financial Economics}, vol.~33, no.~1, pp.~3--56, 1993.
\end{thebibliography}

\end{document}
